\documentclass[10pt,a4paper]{amsart}
\usepackage[margin=19mm]{geometry}
\usepackage{amsmath,amssymb,mathtools}
\usepackage[scr=rsfs,cal=euler]{mathalfa}
\usepackage{xfrac}
\usepackage[unicode,hidelinks,bookmarks=false]{hyperref}
\newcommand{\ie}{i.e.\ }
\newcommand{\cf}{cf.\ }
\newcommand{\resp}{respectively\ }
\newcommand{\Subsection}{\subsection}
\providecommand{\nobreakdash}{\nobreak\hbox{-}}
\setlength{\emergencystretch}{2em}
\multlinegap0pt
\usepackage[all]{xy}
\usepackage{amssymb}
\usepackage{tikz}
\usepackage{float}
\usepackage[shortlabels]{enumitem}
\newtheorem{lemma}{Lemma}
\DeclareMathOperator{\Lip}{Lip}
\title{Positive scalar curvature metrics and aspherical summands}
\author{Shuli Chen, Jianchun Chu,, Jintian Zhu}
\date{Source: arXiv:2312.04698v3 (CC BY 4.0)}
\begin{document}
\maketitle

\noindent\emph{Source and licence.} Positive scalar curvature metrics and aspherical summands. Version arXiv:2312.04698v3 (CC BY 4.0), distributed by arXiv under the Creative Commons Attribution 4.0 International licence, http://creativecommons.org/licenses/by/4.0/, as recorded in the arXiv licence metadata for this exact version and shown on the abstract page https://arxiv.org/abs/2312.04698v3. Full licence notice: \texttt{licenses/arxiv-2312.04698-CC-BY-4.0.txt}.\par\smallskip

\noindent\textit{Editorial scope.} Counted unit: the article's lemma lifting (source0.tex lines 1674--1676). The article's complete original proof of it is delivered with it (source0.tex lines 1677--1781, one \texttt{proof} environment of 105 lines). No mathematics is written, repaired or re-derived: the statement and the proof are the article's own lines, in the article's own notation, and no retained line is substituted or re-flowed.  Retained uncounted as the source-local context the proof needs: lines 1662--1664; lines 1666--1671.  Nothing was omitted from any retained span, so the package carries no omission record.  No retained line contains a citation, so the wrapper carries no bibliography and no bibliography entry.  No retained line contains a figure environment, an image-inclusion command or drawing code, so no image is delivered and the package declares no asset dependency.  Editorial formatting only, in addition: an AMS \texttt{amsart} preamble that restates the article's own declarations and macros used by the retained lines, namely the environments \texttt{lemma} on separate counters, together with the standard packages those lines need.  The retained environments print in this document as Lemma 1 in order of appearance, whereas the article prints the same items with the numbering of its own full document.  The complete native source of the article as distributed in the arXiv e-print for this exact version is bundled unchanged as \texttt{sources/151173/source0.tex}.
\begin{equation}\label{Eq: lifting}
\ker\big(\pi_1(M,x_0)\to \pi_1(N,y_0)\big)=(p_M)_*(\pi_1(\widehat M,\hat x_0)),
\end{equation}

\begin{equation}\label{Eq: lifting diagram}
\begin{split}
\xymatrix{(\widehat M,\hat x_0)\ar[r]^{\hat f}\ar[d]^{p_M}&(\tilde N,\tilde y_0)\ar[d]^{p_N}\\
(M,x_0)\ar[r]^{f}&(N,y_0).}
\end{split}
\end{equation}

\begin{lemma}\label{Lem: lifting}
The map	$\hat f:\widehat M\to \tilde N$ is a smooth quasi-proper map with non-zero degree and we have $|\deg \hat f|\,\big|\,|\deg f|$.
\end{lemma}

\begin{proof}
We first observe that for any $\hat{z}\in\widehat{M}$, the restricted map $\hat f:\mathcal O_{\hat z}\to \mathcal O_{\tilde w}$ is injective, where
\[
\tilde{w} = \hat{f}(\hat{z}), \
z = p_{M}(\hat{z}), \
w = p_{N}(\tilde{w}), \
\mathcal O_{\hat z} = p_M^{-1}(z), \
\mathcal O_{\tilde w} = p_N^{-1}(w).
\]
Otherwise, we can find a curve $\hat\gamma:[0,1] \to \widehat M$ such that $\hat\gamma(0)$ and $\hat\gamma(1)$ are two different points in $\mathcal O_{\hat z}$. Then $p_M\circ \hat\gamma$ induces an element in $\ker\big(\pi_1(M,x_0)\to \pi_1(N,y_0)\big)$ but not in $(p_M)_*(\pi_1(\widehat M,\hat x_0))$, which contradicts \eqref{Eq: lifting}.

Now we show that the set
$$
\tilde S_\infty=\bigcap_{\widehat K\subset \widehat M \text{ compact}} \overline{\hat f(\widehat M- \widehat K)} \subset \tilde N
$$
consists of discrete points. Since $p_N$ is a covering map, $\tilde S_\infty$ consists of discrete points if and only if $p_N(\tilde S_\infty)$ consists of discrete points. It then suffices to prove $p_N(\tilde S_\infty)\subset S_\infty$. For any $\tilde{w}\in\tilde S_\infty$, there exists a sequence $\{\hat{z}_{i}\}\subset\widehat{M}$ such that
\[
\tilde{w} = \lim_{i\to\infty}\hat{f}(\hat{z}_{i}), \ \
\lim_{i\to\infty}d_{g_{\widehat{M}}}(\hat{x}_{0},\hat{z}_{i}) = +\infty,
\]
where $g_{\widehat{M}}$ and $g_{\tilde{N}}$ denote the lifted metrics of $g_{M}$ and $g_{N}$ on $\widehat{M}$ and $\tilde{N}$ respectively. Set $z_{i}=p_{M}(\hat{z}_{i})$. Using the commutative diagram \eqref{Eq: lifting diagram} it is clear that
\[
\lim_{i\to\infty}f(z_{i})
= \lim_{i\to\infty}f(p_{M}(\hat{z}_{i}))
= \lim_{i\to\infty}p_{N}(\hat{f}(\hat{z}_{i}))
= p_{N}(\tilde{w}).
\]
We claim $\lim_{i\to\infty}d_{g_{M}}(x_{0},z_{i})=+\infty$. It then follows that $p_{N}(\tilde{w})\in S_\infty$ as required. To prove the claim, we argue by contradiction. Suppose that the claim is not true. Then after passing to a subsequence, we have $d_{g_{M}}(x_{0},z_{i})\leq L$ for some constant $L$ independent of $i$. Then for each $z_{i}$, there exists $\hat{x}_{i}\in\mathcal{O}_{\hat{x}_{0}}$ such that $d_{g_{\widehat{M}}}(x_{0},z_{i})\leq L$ and so
\[
d_{g_{\tilde{N}}}(\hat{f}(\hat{x}_{i}),\hat{f}(\hat{z}_{i}))
\leq \Lip f|_{B_{L}^{g_{M}}(x_{0})}\cdot d_{g_{\widehat{M}}}(x_{0},z_{i})
\leq \Lip f|_{B_{L}^{g_{M}}(x_{0})}\cdot L.
\]
Set $\tilde{L}=\Lip f|_{B_{L}^{g_{M}}(x_{0})}\cdot L$. Then when $i$ is sufficiently large, we have
\[
d_{g_{\tilde{N}}}(\tilde{w},\hat{f}(\hat{x}_{i}))
\leq d_{g_{\tilde{N}}}(\hat{f}(\hat{x}_{i}),\hat{f}(\hat{z}_{i}))+d_{g_{\tilde{N}}}(\tilde{w},\hat{f}(\hat{z}_{i}))
< \tilde{L}+1.
\]
Since the restricted map $\hat f:\mathcal O_{\hat z}\to \mathcal O_{\tilde w}$ is injective, the set $\mathcal{O}_{\tilde{w}}\cap B_{\tilde{L}+1}^{g_{\tilde{N}}}(\tilde{w})$ is infinite, which is impossible.

Next we want to compute the degree of $\hat f$.
Fix an orientation on $M$ and on $N$ respectively, and then we denote the corresponding fundamental classes by $[M]$ and $[N]$. For any point $x\in M$ we shall let $[M]_x$ denote the image of $[M]$ under the map $H_n^{\mathrm{lf}}(M)\to H_n(M,M-\{x\})$. As a convention, similar notations will be used later without further explanation when there is no confusion caused. As a natural choice, we take the orientation on $\widehat M$ and $\tilde N$ such that the corresponding fundamental classes $[\widehat M]$ and $[\tilde N]$ satisfy
$$
(p_M)_{*}([\widehat M]_{\hat x})=[M]_{p_M(\hat x)} \mbox{ for all }\hat x\in \widehat M
$$
and
$$
(p_N)_{*}([\tilde N]_{\tilde y})=[N]_{p_N(\tilde y)} \mbox{ for all }\tilde y\in \tilde N.
$$

Let us consider the restricted map of $f$ on $f^{-1}(N-S_\infty)$ denoted by
$$
f_{\mathrm{res}}:f^{-1}(N-S_\infty)\to N-S_\infty.
$$
Recall that the map $f:M\to N$ has non-zero degree and so the map $f_{\mathrm{res}}$
is surjective. Based on a change of marked points we may assume that $y_0 $ is a point contained in $ N-S_\infty$ and that it is also a regular value of the map $f$.

Let us divide our discussion into two cases:

\medskip
{\it Case 1. the map $f_*:\pi_1(M,x_0)\to \pi_1(N,y_0)$ is surjective.} Notice that the map $f_{\mathrm{res}}$ is proper, so the preimage $f^{-1}(y_0)$ consists of finitely many (regular) points and we can do labeling
$$
f^{-1}(y_0)=\{x_0,x_1,\ldots,x_l\}.
$$
Then we have
$$
\sum_{i=0}^lf_*([M]_{x_i})=\deg f\cdot [N]_{y_0}.
$$
For each $i$ let us fix a point $\hat x_i\in\widehat M$ satisfying $p_M(\hat x_i)=x_i$.

We have already known that the restricted map $\hat f:\mathcal O_{\hat x_i}\to \mathcal O_{\tilde y_0}$ is injective. We show it is also surjective. To see this we denote $\tilde y_i=\hat f(\hat x_i)$ and for any point $\tilde y_0'\in\mathcal O_{\tilde{y}_0}$ we take a curve $\tilde \gamma:[0,1]\to \tilde N$ connecting $\tilde y_i$ and $\tilde y'_0$. Correspondingly we obtain a closed curve $\gamma:(\mathbb S^1,1)\to (N,y_0)$. Recall that the map $f_*: \pi_1(M,x_0)\to \pi_1(N,y_0)$ is surjective by assumption and the same thing holds with $\pi_1(M,x_0)$ replaced by $\pi_1(M,x_i)$. As a result, we can find a closed curve $\gamma_i:(\mathbb S^1,1)\to (M,x_i)$ such that $f_*([\gamma_i])=[\gamma]$. By lifting we can obtain a curve $\hat\gamma_i:[0,1]\to \widehat M$ with $\hat\gamma_i(0)=\hat x_i$ and $p_M\circ \hat\gamma_i=\gamma_i$. Then it is easy to verify that $\hat\gamma_i(1)$ is a point in $\mathcal O_{\hat x_i}$ with $\hat f$-image $\tilde y_0'$.

Now we can compute
\[
\begin{split}
\deg \hat f\cdot [\tilde N]_{\tilde y_0}&=\sum_{\hat x\in \hat f^{-1}(\tilde y_0)}\hat f_*([\widehat M]_{\hat x})\\
&=\sum_{i=0}^l\sum_{\hat x\in \hat f^{-1}(\tilde y_0)\cap \mathcal O_{\hat x_i}}\hat f_*([\widehat M]_{\hat x}).
\end{split}
\]
The above discussion shows that $\hat{f}^{-1}(\tilde{y}_{0})\cap\mathcal{O}_{\hat{x}_{i}}$ consists of one point. Applying $(p_{N})_{*}$ on both sides, we obtain
$$
\deg \hat f\cdot [N]_{y_0}=\sum_{i=0}^l f_*([M]_{x_i})=\deg f\cdot [N]_{y_0}.
$$
Therefore we obtain $\deg \hat f=\deg f$ and in particular $|\deg\hat f|\,\big |\,|\deg f|$.

\medskip
{\it Case 2. the map $f_*:\pi_1(M,x_0)\to \pi_1(N,y_0)$ is not surjective.} In this case, we consider the covering $\check p_N:(\check N, \check y_0) \to (N,y_0)$ such that
$$
(\check p_N)_*(\pi_1(\check N,\check y_0))=f_*(\pi_1(M,x_0))\subset \pi_1(N,y_0).
$$
In particular, we can lift the map $f$ to a map $\check f:(M,x_0)\to (\check N,\check y_0)$ such that $\check{f}_{*}:\pi_{1}(M,x_{0})\to\pi_{1}(\check{N},\check{y}_{0})$ is surjective. From the discussion in Case 1 we conclude $\deg \check f=\deg \hat f$. Let us label
$$
\check p_N^{-1}(y_0)=\{\check y_j\}_{j\in J}.
$$
Notice that we have the relation
\[
\begin{split}
\deg f\cdot [N]_{y_0}&=\sum_{i=0}^l f_*([M]_{x_i})\\
&=\sum_{j\in J}\sum_{x_i\in \check f^{-1}(\check y_j)}{(\check{p}_N)_*}\check f_*([M]_{x_i})\\
&=|J|\cdot \deg \check f\cdot [N]_{y_0},
\end{split}
\]
which implies $|J|$ is finite and so $|\deg\hat f|\,\big |\,|\deg f|$.
\end{proof}

\end{document}
